\documentclass[11pt]{article}

% Camera-ready ACL style (use [review] for the anonymised submission version).
\usepackage{acl}
\usepackage{times}
\usepackage{latexsym}
\usepackage{amsmath}
\usepackage[T1]{fontenc}
\usepackage[utf8]{inputenc}
\usepackage{microtype}
\usepackage{booktabs}
\usepackage{graphicx}
\usepackage{multirow}
\usepackage{enumitem}

\title{STARK-QA: A Reproducible Hybrid RAG--KAG System for\\Narrative Question Answering on FairytaleQA}

\author{Ishaan Romil (2023114011) \quad Archit Choudhary (2023114002) \quad
  Sujal Deoda (2022115001) \\
  International Institute of Information Technology, Hyderabad}

\begin{document}
\maketitle

\begin{abstract}
We present \textbf{STARK-QA} (\emph{Story-grounded Text Answering with Retrieval
and Knowledge}), a question-answering system for the FairytaleQA narrative
comprehension benchmark that fuses dense\,+\,sparse passage retrieval (RAG) with
a lightweight knowledge graph (KAG) and generates grounded answers with a
large language model. Beyond the system itself, this report documents a
\emph{reproducibility and repair} effort: we took a non-running research
prototype, restored it to a fully working state, hardened every external
dependency with a graceful fallback, and pinned a locked, one-command
environment. On 100 held-out FairytaleQA test questions, STARK-QA attains a
token-level F1 of \textbf{0.611}, \textbf{3.5$\times$} the Basic-RAG ablation
(0.176) and \textbf{14.9$\times$} the Basic-KAG ablation (0.041), while leading
on Exact Match, ROUGE-1/2/L, BLEU, semantic similarity and BERTScore. We release
the full pipeline, cached data splits, and a deterministic \texttt{uv} lockfile.
\end{abstract}

\section{Introduction}
Narrative question answering requires a system to read a story and answer
free-form questions about its characters, settings, actions, and the causal and
emotional relationships between them. FairytaleQA \citep{xu-etal-2022-fairytaleqa}
is an authentic benchmark of such questions written by education experts over
classic fairy tales, spanning seven attribute types (\emph{character},
\emph{setting}, \emph{action}, \emph{feeling}, \emph{causal relationship},
\emph{outcome resolution}, and \emph{prediction}) and an explicit/implicit
reasoning split.

Two paradigms dominate knowledge-intensive QA. Retrieval-augmented generation
(RAG) \citep{lewis2020rag} retrieves relevant passages and conditions a generator
on them; it excels at descriptive \emph{where}/\emph{why} questions whose answer
is a span of text. Knowledge-augmented generation (KAG) instead extracts
structured entities and relations \citep{pan2024kgsurvey} and is naturally suited
to entity-centric \emph{who} questions. STARK-QA \textbf{fuses both}: it orders
RAG and KAG evidence by the question type and lets an LLM compose a short,
grounded answer, so it inherits the strengths of each rather than committing to
one.

This report makes three contributions:
\begin{enumerate}[noitemsep]
    \item A hybrid RAG\,+\,KAG architecture with question-type-aware evidence
    fusion and grounded LLM generation (\S\ref{sec:system}).
    \item A reproducibility/repair effort that turns a broken prototype into a
    one-command, dependency-locked, gracefully degrading pipeline
    (\S\ref{sec:repro}).
    \item A full evaluation against RAG-only and KAG-only ablations across eight
    metrics and seven question types (\S\ref{sec:results}).
\end{enumerate}

\section{System Architecture}
\label{sec:system}
Given a question $q$ and the story corpus reconstructed from the test split,
STARK-QA produces an answer in four stages.

\paragraph{Hybrid retrieval (RAG).} Stories are chunked into overlapping
word windows (200 words, overlap 40). We embed chunks with a
Sentence-BERT encoder \citep{reimers-gurevych-2019-sentencebert}
(\texttt{all-MiniLM-L6-v2}) and index them with FAISS
\citep{johnson2021faiss} for dense semantic search, run in parallel with a
TF--IDF lexical index. The two score vectors are min--max normalised and fused as
$s = \alpha\, s_{\text{dense}} + (1-\alpha)\, s_{\text{tfidf}}$ with
$\alpha = 0.6$, followed by a light question-type re-ranking pass.

\paragraph{Knowledge graph (KAG).} From each story we extract entities
(characters, objects/places, descriptive concepts), simple co-occurrence
relationships, and the sentences in which each entity appears, all with
dependency-free regular expressions. At query time the question is matched to
entities (exact $\rightarrow$ partial $\rightarrow$ semantic) and the most
relevant stored sentences plus matching relationships are returned.

\paragraph{Question-type-aware fusion.} RAG and KAG evidence are ordered by the
kind of question---entity-first for \emph{who}, passage-first for
\emph{where}/\emph{why}---before generation, so the prompt foregrounds the most
useful evidence for that type.

\paragraph{Grounded generation.} A few-shot, FairytaleQA-styled prompt conditions
the Mistral LLM (\texttt{mistral-small-latest}, temperature $0$) on the fused
evidence to produce the shortest complete answer span. When no API key is
present, the system falls back to a dependency-free extractive answerer over the
same fused evidence, so it always returns a grounded answer.

\paragraph{Ablation baselines.} \textbf{Basic-RAG} uses TF--IDF retrieval with
purely extractive answering (no dense embeddings, no KAG). \textbf{Basic-KAG}
builds an entity co-occurrence graph and answers from graph neighbourhoods. Both
are deliberately weaker, isolating the contribution of hybrid retrieval and of
the fusion-plus-LLM design.

\section{Reproducibility and Engineering}
\label{sec:repro}
A substantial part of this work was restoring a non-running prototype to a
correct, reproducible state. The salient changes were:
\begin{itemize}[noitemsep]
    \item \textbf{Shared utilities.} Sentence splitting, tokenisation, answer
    normalisation and corpus fingerprinting were centralised, removing
    behavioural drift between components that previously defined these
    differently.
    \item \textbf{Graceful degradation.} Every heavy dependency
    (\texttt{torch}/\texttt{sentence-transformers}/FAISS, spaCy, NLTK, ROUGE,
    BERTScore, Matplotlib) has a pure-Python or TF--IDF fallback, so a missing
    package degrades quality rather than crashing the run.
    \item \textbf{LLM client compatibility.} Generation now probes the modern,
    namespaced and legacy Mistral SDK import layouts and uses a current model id,
    fixing breakage against newer SDK releases; calls are throttled and retried
    with exponential backoff to respect rate limits.
    \item \textbf{Offline data and caching.} The data loader imports
    \texttt{datasets} lazily and reuses cached JSON splits offline; retrieval
    indices are cached and keyed by a fingerprint of the corpus so changing the
    data transparently rebuilds the index.
    \item \textbf{Locked environment.} Dependencies are declared in
    \texttt{pyproject.toml} and pinned in a \texttt{uv.lock} (107 packages), so
    \texttt{uv sync} reproduces the exact stack on Python~3.11.
\end{itemize}

\section{Experimental Setup}
\label{sec:setup}
\paragraph{Data.} We evaluate on the first 100 questions of the FairytaleQA test
split, shipped as a cached JSON split for fully offline reproduction. The primary
human reference answer is used as ground truth.

\paragraph{Metrics.} We report Exact Match and token-level F1 (SQuAD-style
normalisation), ROUGE-1/2/L \citep{lin-2004-rouge}, BLEU
\citep{papineni-2002-bleu}, embedding-based semantic similarity, and BERTScore
F1 \citep{zhang2020bertscore}, together with mean response time. BERTScore uses
\texttt{distilbert-base-uncased} to avoid a large model download; the canonical
\texttt{roberta-large} scorer is configurable.

\paragraph{Configuration.} Retrieval top-$k=5$, hybrid weight $\alpha=0.6$, chunk
size 200 words. Generation uses temperature $0$ and an 80-token cap. All settings
live in \texttt{config.py} and are overridable via environment variables.

\section{Results}
\label{sec:results}
Table~\ref{tab:main} reports the headline comparison. STARK-QA leads on
\emph{every} metric. Its token-F1 of 0.611 is $3.5\times$ Basic-RAG and
$14.9\times$ Basic-KAG, and the gap widens on the generation-oriented metrics
(BLEU $0.422$ vs.\ $0.085$; ROUGE-2 $0.467$ vs.\ $0.128$), reflecting the value
of grounded LLM composition over purely extractive answering. The cost is
latency: STARK-QA averages $1.73$\,s per question (an API round-trip plus
proactive throttling) against sub-millisecond local baselines.

\begin{table}[t]
\centering
\small
\begin{tabular}{@{}lccc@{}}
\toprule
\textbf{Metric} & \textbf{STARK-QA} & \textbf{Basic-RAG} & \textbf{Basic-KAG} \\
\midrule
Exact Match    & \textbf{0.35}  & 0.00 & 0.01 \\
Token F1       & \textbf{0.611} & 0.176 & 0.041 \\
ROUGE-1        & \textbf{0.629} & 0.196 & 0.058 \\
ROUGE-2        & \textbf{0.467} & 0.128 & 0.010 \\
ROUGE-L        & \textbf{0.627} & 0.186 & 0.055 \\
BLEU           & \textbf{0.422} & 0.085 & 0.008 \\
Semantic sim.  & \textbf{0.729} & 0.381 & 0.193 \\
BERTScore F1   & \textbf{0.893} & 0.729 & 0.657 \\
\midrule
Avg.\ time (s) & 1.726 & 0.001 & 0.003 \\
\bottomrule
\end{tabular}
\caption{Main results on 100 FairytaleQA test questions. Baselines are
deterministic; STARK-QA figures are from the reported Mistral run.}
\label{tab:main}
\end{table}

\paragraph{Per-type analysis.} Table~\ref{tab:bytype} breaks F1 down by question
attribute. STARK-QA is strongest on \emph{feeling} (1.00) and \emph{character}
(0.82) questions, where the fused entity and emotion evidence is highly
predictive, and weakest on \emph{prediction} (0.36) and \emph{outcome
resolution} (0.35), which require inference beyond the literal text. The ordering
is consistent across systems---every type that is hard for the baselines is also
the relative weak point for STARK-QA---but STARK-QA dominates within each type,
e.g.\ \emph{feeling} improves from $0.03$/$0.00$ (RAG/KAG) to $1.00$.

\begin{table}[t]
\centering
\small
\begin{tabular}{@{}lccc@{}}
\toprule
\textbf{Question type} (n) & \textbf{STARK} & \textbf{B-RAG} & \textbf{B-KAG} \\
\midrule
feeling (8)              & \textbf{1.00} & 0.03 & 0.00 \\
character (9)            & \textbf{0.82} & 0.10 & 0.11 \\
causal relationship (19) & \textbf{0.64} & 0.26 & 0.08 \\
action (37)             & \textbf{0.58} & 0.20 & 0.01 \\
setting (11)            & \textbf{0.58} & 0.10 & 0.03 \\
prediction (7)          & \textbf{0.36} & 0.13 & 0.01 \\
outcome resolution (9)  & \textbf{0.35} & 0.24 & 0.09 \\
\bottomrule
\end{tabular}
\caption{Token-F1 by FairytaleQA question type. STARK-QA dominates within every
type; the relative ordering of difficulty is shared across systems.}
\label{tab:bytype}
\end{table}

\paragraph{Discussion.} The ablations confirm the design hypothesis. Basic-RAG
(lexical-only, extractive) recovers some content on descriptive
\emph{causal}/\emph{action} questions but cannot phrase a clean short answer,
capping its F1 near $0.18$ and its Exact Match at $0$. Basic-KAG is weaker still:
co-occurrence neighbourhoods rarely contain the precise answer span. STARK-QA's
combination of dense+sparse retrieval, structured KAG context, and LLM
composition lifts both the \emph{retrieval} of relevant evidence and the
\emph{surface form} of the answer, which is why the gain is largest on the
generation-sensitive metrics (BLEU, ROUGE-2) and on Exact Match.

\section{Conclusion}
STARK-QA shows that fusing hybrid retrieval with a lightweight knowledge graph
and grounding an LLM on the result yields large, consistent gains over either
component alone on narrative QA. Equally important for a course/research artifact,
the system is now \emph{reproducible}: a single \texttt{uv sync} reconstructs the
locked environment, cached splits enable offline runs, and graceful fallbacks
keep the pipeline running even when optional dependencies are absent.

\section*{Limitations}
The evaluation covers 100 test questions from a single split; broader coverage
and multiple reference answers would tighten the estimates. STARK-QA's
generation depends on a proprietary LLM, so its scores are not perfectly
reproducible run-to-run (the deterministic extractive fallback trades some
quality for full reproducibility). Entity and relation extraction in the KAG
module is regex-based and English/fairy-tale-specific, and BERTScore uses a
lighter encoder than the canonical \texttt{roberta-large}, which may shift
absolute BERTScore values.

\section*{Reproducibility}
Code, cached FairytaleQA splits, configuration and the \texttt{uv.lock} are
included in the repository. To reproduce:
\texttt{uv sync -{}-python 3.11}; \texttt{uv run python -m spacy download
en\_core\_web\_sm}; \texttt{uv run python main.py -{}-test\_size 100}. Outputs are
written to \texttt{results/}.

\bibliography{custom}

\end{document}
